\section{Materials and Methods}
\label{sec:methods}

% =====================================================
\subsection{Data sources}
\label{sec:data}

The task is relation extraction at the abstract level: given a PubMed
abstract and two pre-identified entity spans, predict the relation
type from a schema vocabulary (or \code{\_\_NEGATIVE\_\_} if no
relation holds). We trained on three public corpora: BioRED
\citep{luo2022biored} (600 abstracts, multiple relation types across
gene, disease, variant, and chemical entities; 500/100
train-dev/test split); DrugProt \citep{miranda2023drugprot}
(4{,}250 documents, thirteen chemical--protein mechanism relations);
and BC5CDR \citep{li2016bc5cdr} (1{,}000 train+dev + 500 test
documents, binary chemical-induces-disease relations).

The knowledge-base anchor was CIViC \citep{griffith2017civic}. We
extracted 165 evidence-supported targets (PubMed abstract plus two
entity spans). Of these, 162 mapped to schema vocabularies under
the family-to-label projection described in \S\ref{sec:kb-audit}
and form the evaluable target pool; the remaining three carry
variant--gene endpoints not covered by any of the three schema
vocabularies and are documented as a structural coverage gap. An
oncology-projected training subset comprised 733 abstracts whose
PubMed records carry a Medical Subject Headings (MeSH) C04
(Neoplasms) descriptor \citep{lipscomb2000mesh}. Corpus details are
in Supplement~B.

% =====================================================
\subsection{Schema design}
\label{sec:schema}

Three nested relation schemas were defined at increasing granularity.
The 4-label \shead{family} schema (S\textsubscript{family}) is the
minimum partition preserving three clinical pair-type categories:
drug--disease, drug--gene regulation, and other associations.%
\footnote{S\textsubscript{family} is named \texttt{S\_flat} in the
pre-registration document and analysis code; we rename it here for
readability.} The
8-label \shead{pair} schema (S\textsubscript{pair}) decomposes the general-association category into
entity-pair-specific categories (\code{GENE\_DISEASE},
\code{VARIANT\_DISEASE}, etc.), recovering the entity-pair
distinctions that BioRED natively annotates. The 13-label
\shead{mechanism} schema (S\textsubscript{mech}) further decomposes
drug--gene regulation into five DrugProt sub-mechanisms. A relation is eligible for training and evaluation only if its
entity-pair type is covered by the schema. Schema vocabularies,
legal-pair filters, and the projection used in the KB audit are in
Supplement~C.

% =====================================================
\subsection{Training protocol}
\label{sec:training}

Four encoders were trained: PubMedBERT-base
\citep{gu2021pubmedbert}, BioLinkBERT-base
\citep{yasunaga2022biolinkbert}, PubMedBERT-large
\citep{gu2021pubmedbert}, and a general-domain reference
(RoBERTa-base \citep{liu2019roberta}). Three training schedules were
compared: single-corpus (BioRED only); multi-corpus joint (BioRED +
DrugProt + BC5CDR in one stage); and oncology-projected staged, in
which a continuation stage on the MeSH-C04 subset followed the
multi-corpus stage. Each training stage used a fixed budget of
2{,}048 optimiser updates with a cap of 2{,}000 positive rows per
corpus shard. The staged schedule consumed 4{,}096 updates in
total, so pairwise contrasts within an encoder compare schedule
composition rather than total compute. All runs share the same
optimiser, learning rate, and early-stopping criterion (full
hyperparameter inventory in Supplement~D).

The study comprised two waves. The \emph{schema-comparison wave}
($n = 120$ runs) varied schema and encoder under a fixed schedule
to select the operational schema (RQ1). The \emph{configuration
wave} ($n = 180$ main runs plus a 10-seed RoBERTa-base reference
cell, 190 total) fixed schema to the first-wave selection and
varied encoder $\times$ schedule with 20 seeds per cell,
addressing RQ2, RQ3, and RQ4. Seed indices 1--10 are shared across
the two waves so that cross-wave paired contrasts can match;
indices 11--20 are used in the configuration wave only.
Hyperparameters and seeding policy are in Supplement~D.

A pre-registered parameter-efficient fine-tuning arm based on
low-rank adaptation \citep[LoRA;][]{hu2022lora} (rank 16; attention
query and value projections; classifier head fully trained)
collapsed to uniform negative prediction under all three
pre-committed configurations and was declared a methodological null
(Supplement~I).

% =====================================================
\subsection{Knowledge-base audit}
\label{sec:kb-audit}

Each of the 162 CIViC targets carries a PubMed abstract, two
pre-identified entity spans, and a curator-recorded relation that
maps to one or more schema-vocabulary labels under a heuristic
family-to-label projection. Because entity spans are pre-identified,
the audit isolates relation-type discrimination rather than entity
recognition. The projection maps each (entity-pair family, CIViC
evidence-type, schema) tuple to an expected-label set, which is
set-valued when the CIViC evidence type is compatible with more
than one schema label; the full projection rules and the released
JSON are in Supplement~C.

\paragraph{Inter-annotator audit.}
We audited the heuristic projection against an independent second
annotator on a stratified random sample of 30 of the 162 evaluable
targets (18.5\%). The second annotator was a large language model
(Claude Opus 4.7, \texttt{temperature\,=\,0}, deterministic
decoding) prompted with the same label definitions but blind to the
heuristic's output \citep{gilardi2023llmannotator}. The resulting
$\kappa$ is a model-based validity check rather than a
human--human inter-annotator estimate, and is reported as a lower
bound on agreement with an expert human curator. Observed agreement
was \(\kappa = 0.56\) (95\% CI [0.32, 0.80]); all seven
disagreements followed a single directional pattern detailed in
\S\ref{sec:results-iaa}. Per-target rationales are in Supplement~C.

\paragraph{KB-audit metrics.}
For each model on each target we computed three correctness-aware
metrics: (i)~\textbf{KB argmax accuracy} (primary): 1 if the
model's top-ranked label is in the expected-label set, 0 otherwise;
(ii)~\textbf{KB truth probability}: the model's predicted
probability summed over labels in the expected set; and
(iii)~\textbf{KB selective-recall AUC}: the area under the curve of
per-target precision against rejection fraction, computed by
sweeping a confidence threshold and abstaining on predictions below
it. Intuitively, (i) asks whether the model's top prediction is
right, (ii) asks how much probability mass landed on truth, and
(iii) asks how reliably the model abstains when wrong. KB argmax
accuracy was the pre-registered primary metric; KB truth probability
and KB selective-recall AUC are reported as discrimination and
selective-prediction sensitivity analyses (\S\ref{sec:results-rq3}).
Formal definitions are in Supplement~C.

\paragraph{External benchmarks.}
The primary external benchmark is the macro-averaged F1 score
(macro-F1) of BioRED over its non-negative labels (BioRED ex-NEG);
BC5CDR drug--disease F1 is the out-of-domain benchmark.

% =====================================================
\subsection{Statistical analysis}
\label{sec:stats}

\paragraph{Target estimands.}
Each contrast estimates the observed effect of one design lever
across the realised factorial cells, with seeds treated as
nuisance variation; we do not extrapolate variance shares to a
wider population of encoders or schedules.

\paragraph{Reliability of the primary KB metric.}
KB argmax accuracy is noisier across seeds than the external
benchmarks. On the configuration wave (9 cells $\times$ 20 seeds),
the intraclass correlation (one-way random effects; ICC(1,1)
\citep{shrout1979intraclass}) for KB argmax accuracy was 0.67,
indicating substantial between-cell separation, but the pooled
within-cell standard deviation was 0.17 --- about three times the
within-cell standard deviation for BioRED ex-NEG F1 over the same
nine cells (0.06). Pre-registered seed-paired contrasts
absorb cell-level noise into the within-cell difference and were
specified for every KB-side comparison; the residual seed-level
dispersion still widens paired-bootstrap intervals at the realised
cell counts and is the precision constraint behind reporting the
coupling-slope family descriptively (\S\ref{sec:results-rq4}).

\paragraph{Comparison protocol.}
The locked primary family contained 21 planned comparisons
(Supplement~K). After H4 (LoRA) was declared a methodological null
and H5 (architecture) was deferred, the realised analysis covers 15
comparisons: nine paired-difference contrasts entered into the
Benjamini--Hochberg family (H1$\,$two encoder-scale contrasts;
H2$\,$one multi-corpus contrast; H3$\,$six staging contrasts), five
coupling slopes (H6, reported descriptively), and one variance
ratio (H7). The paired-difference contrasts use
Benjamini--Hochberg false-discovery-rate adjustment
\citep{benjamini1995controlling} at \(q=0.05\); each reports both
paired \(t\) and Wilcoxon signed-rank
\citep{wilcoxon1945individual}, with paired \(t\) as headline.
Paired-bootstrap 95\% confidence intervals (\(B=10{,}000\)) resample
within-cell differences across seeds and form the primary
uncertainty estimate.

\paragraph{Variance decomposition.}
For each wave we partitioned the sum-of-squares of BioRED ex-NEG F1
and of KB argmax accuracy across design factors using a Type-I
analysis of variance (sequential fits in the pre-registered factor
order; Supplement~K). The variance ratio \(R\) is the ratio of the
design-lever variance share in the benchmark metric to the same
share in the KB metric: \(R<1\) means the same design levers
explain more KB variance than benchmark variance, and \(R>1\)
inverts that direction. \(R_B\) is accompanied by a
cluster-bootstrap 95\% CI (5{,}000 cell-level resamples). We also
report the rank-inversion rate: the fraction of configuration-wave
cell pairs whose BioRED ex-NEG F1 differs by less than the
pre-registered matching radius \(\rho=0.03\) (one schema-wave
within-cell BioRED SD) and whose KB rankings disagree. A sensitivity
sweep over \(\rho\in\{0.01, 0.03, 0.05\}\) is in Supplement~H.

% =====================================================
\subsection{Pre-registration and protocol deviations}
\label{sec:prereg}

The factorial design, hypotheses, decision rules, and
multiple-comparison tiering were pre-registered
\citep{nosek2018preregistration} before any configuration-wave run.
The pre-registration document and amendment log are in the Data and
Code Availability statement.

Three post-lock amendments were documented in the versioned amendment
log. \emph{Amendment~1} (pre-unblinding): a slope-CI reportability
gate (\(\leq 0.30\) width) was added after the schema-comparison
wave revealed wider within-cell dispersion than assumed; the
variance-ratio factor set was also clarified.
\emph{Amendment~2}: the LoRA hypothesis (H4) was declared null
after all configurations collapsed; the architecture hypothesis
(H5) was deferred. During manuscript preparation we also
reclassified the PubMedBERT-base vs BioLinkBERT-base contrast,
originally listed under H1 encoder-scale, as an
architecture comparison rather than a scale comparison (both
encoders are base-scale; they differ in pretraining corpus and
initialisation). It is reported in Supplement~F alongside the
remaining H1 results but is not part of the encoder-scale family
in Table~\ref{tab:t2}, consistent with H5's deferred status.
\emph{Amendment~3}: the inter-annotator audit, the ICC analysis,
and the \(\rho\) sensitivity analysis were added as post-hoc
validity checks during manuscript preparation. No pre-registered
decision rule, threshold, or comparison tier was amended after
unblinding.

% =====================================================
\subsection{Computational environment}

Training and evaluation ran on the University of Bristol's BriCS
Isambard-AI system under Slurm; the pinned software environment is
in Supplement~A.
